\clearpage
\VerdierPhysicalPageAnchor{0239}%
\begingroup\small\setlength{\parskip}{0pt}\setlength{\emergencystretch}{1.5em}
\setlength{\abovedisplayskip}{4pt plus 2pt minus 2pt}%
\setlength{\belowdisplayskip}{4pt plus 2pt minus 2pt}%

\VerdierSemanticAnchor{ch4d-4-4-2-b-cont}
\noindent\textit{于是有：}
\[
\mathrm I_2^{p,q}=H_{\mathcal B}^pF^{q,\bullet}(Y^\bullet).
\]

\VerdierSemanticAnchor{ch4d-4-4-2-c}
\noindent\textit{c) 各项 $\mathrm I_\infty^{p,q}$ 是适当滤过的对象
$F^{p+q}(Y)$ 的关联分次对象。}

\VerdierSemanticAnchor{ch4d-4-4-2-proof}
证明 a)。设 $\bigl(Y(p,q),\delta(p,q,r)\bigr)$ 是与复形 $Y$ 相联系的、
取值于 $K(\mathcal A)$ 的第一谱对象。由该对象的定义立即得到
(4.1.2.3)：
\[
Y(-p-1,-p)=Y^p[-p].
\]
因此（第二章 4.3.3.1）：
\[
\mathrm I_1^{p,q}=F^{p+q}\bigl(Y^p[-p]\bigr)=F^q(Y^p).
\]
微分 $d_1\colon\mathrm I_1^{p,q}\longrightarrow\mathrm I_1^{p+1,q}$ 的描述
随即由 (4.1.2.6) 直接得出。断言 b) 立即由 a) 得出。最后一个断言由
（第二章 4.3.2）得出。

\VerdierSemanticAnchor{ch4d-4-4-3}
\medskip\noindent\textbf{注 4.4.3.}\quad
用 $\bigl(\mathrm I_r^{p,q}(Y),d_r,r\geq1\bigr)$ 表示函子 $F$ 相对于复形
$Y$ 的第一谱序列。函子
\[
Y\longmapsto\bigl(\mathrm I_r^{p,q}(Y),d_r,r\geq1\bigr)
\]
是从范畴 $\operatorname{comp}(\mathcal A)$ 到取值于 $\mathcal B$ 的谱序列
范畴的函子。一般而言，此函子不能通过范畴 $K(\mathcal A)$ 分解。另一方面，
函子

\VerdierSemanticAnchor{ch4d-4-4-3-r-ge-2}
\[
Y\longmapsto\bigl(\mathrm I_r^{p,q}(Y),d_r,r\geq2\bigr)
\]
由命题 (4.4.2) b)，可通过复形同伦范畴分解。

\VerdierSemanticAnchor{ch4d-4-4-4}
\medskip\noindent\textbf{4.4.4.}\quad
现在设 $\mathcal A$ 是阿贝尔范畴，
$F\colon D(\mathcal A)\longrightarrow\mathcal B$ 是取值于阿贝尔范畴
$\mathcal B$ 的上同调函子，并设 $F^p$，$p\in\mathbb Z$，是由 $F$ 与
$D(\mathcal A)$ 的平移复合所得的函子（$F^p(X)=F(X[p])$）。此外，设 $Y$
是 $\operatorname{comp}(\mathcal A)$ 的对象，并设
$\bigl(Y(p,q),\delta(p,q,r)\bigr)$ 是与复形 $Y$ 相联系的、取值于
$D(\mathcal A)$ 的第一谱对象。函子 $F$ 把该谱对象变为取值于
$\mathcal B$ 的谱对象。由此得到一个取值于 $\mathcal B$ 的谱序列，仍称为
\emph{函子 $F$ 相对于复形 $Y$ 的第一谱序列}。这个谱序列在本质上由命题
(4.4.2) 描述（只须在命题 (4.4.2) 中把范畴 $K(\mathcal A)$ 换成范畴
$D(\mathcal A)$）。函子 $F$ 相对于复形 $Y$ 的第一谱序列，对 $F$ 以及
当论元 $Y$ 在 $\operatorname{comp}(\mathcal A)$ 中变化时的该论元都是函子性的。
对此谱序列，我们始终采用（第二章 4.3.3.1）所引入的指标方式。
\endgroup
